\chapter{凸集与超平面分离}

\label{chap:04}

从第 1 章到第 3 章，我们已经分别准备了拓扑、测度与泛函分析的基础语言，但无穷维凸优化真正的几何核心，直到本章才开始出现。有限维教材中常见的图像是：凸集被超平面支撑，最优性条件由法向量表达，极值点则集中在边界最尖锐的部分。这些直觉在 $\mathbb R^n$ 中几乎不需要额外解释；然而一旦进入函数空间、弱拓扑和一般 Banach 空间，它们就必须被重新翻译成线性泛函、分离定理和极点结构，否则所谓“几何直观”只会停留在图片层面。

本章默认主要工作在实线性空间中，因为凸组合、超平面分离和支撑泛函的自然语言本来就是实的。第 \ref{sec:4-1} 节先把凸集、仿射集、锥、闭包与凸包这些对象重新整理出来，并指出无穷维中“闭凸集没有内点”并不是例外，而是常态。第 \ref{sec:4-2} 节进一步引入代数内部与相对内部，解释为什么第 5 章与第 7 章中的资格条件不能只依赖通常内点。第 \ref{sec:4-3} 节是本章的核心：Hahn--Banach 定理把线性泛函的延拓与几何分离统一到同一条证明链中，从而为次微分、对偶和 KKT 条件提供最根本的泛函来源。最后第 \ref{sec:4-4} 节以极点与 Krein--Milman 定理收束，说明一个紧凸集如何由其边界上的极端结构控制。

从 Boyd 的有限维叙事看，本章做的事情可以概括为一句话：把“超平面、相对内部、支撑和顶点”的欧氏语言提升为无穷维可复用的泛函语言。后文中出现的次梯度、共轭函数、强对偶、广义 KKT 以及原子表示，都要回链到本章的结论。因此阅读这一章时，最重要的不是背诵所有定义，而是看清楚哪一个结论依赖拓扑、哪一个结论依赖局部凸性、哪一个结论又必须退回 Banach 或 Hilbert 环境才真正可用。

\section{无穷维向量空间中的凸集}

\label{sec:4-1}

凸性首先是一种代数稳定性：它要求集合在两点之间的线段上闭合，而并不预设任何坐标、内积或维数。因此，凸集理论天然适用于函数空间、测度空间和算子空间。真正进入无穷维以后，困难并不在“凸”的定义本身，而在于拓扑直觉不再自动成立。有限维里常见的事实，例如闭有界集有较好几何边界、闭凸集常常带有非空内部、凸包与闭包的运算容易可视化，在一般 Banach 或 Hilbert 空间中都需要重新说明。特别是当我们谈论闭包时，实际上已经回到了第 \ref{sec:1-1} 节命题~\ref{prop:closure-net} 所建立的网语言；而当我们谈论最小二乘、正则化与投影时，又必须回到第 \ref{sec:3-1} 节定义~\ref{def:hilbert-space} 所代表的几何环境。

\subsection{基本对象}

\begin{definition}[凸集]\label{def:convex-set}
设 $X$ 为实向量空间，$C\subset X$。若对任意 $x,y\in C$ 与任意 $\lambda\in[0,1]$ 都有
\[
\lambda x+(1-\lambda)y\in C,
\]
则称 $C$ 为一个\emph{凸集}。
\end{definition}

\begin{definition}[仿射集]\label{def:affine-set}
设 $X$ 为实向量空间，$A\subset X$。若对任意 $x,y\in A$ 与任意 $\lambda\in\mathbb R$ 都有
\[
\lambda x+(1-\lambda)y\in A,
\]
则称 $A$ 为一个\emph{仿射集}。等价地，$A$ 是某个线性子空间的平移。
\end{definition}

\begin{definition}[凸锥]\label{def:convex-cone}
设 $X$ 为实向量空间，$K\subset X$。若对任意 $x,y\in K$ 与任意 $\alpha,\beta\ge 0$ 都有
\[
\alpha x+\beta y\in K,
\]
则称 $K$ 为一个\emph{凸锥}。
\end{definition}

若 $E\subset X$，则记
\[
\operatorname{conv}(E)
\]
为包含 $E$ 的最小凸集，称为 $E$ 的\emph{凸包}；记
\[
\operatorname{aff}(E)
\]
为包含 $E$ 的最小仿射集，称为 $E$ 的\emph{仿射包}；记
\[
\operatorname{cone}(E)
\]
为包含 $E$ 的最小凸锥。只要 $X$ 还带有拓扑，我们就继续使用 $\overline{E}$ 和 $\operatorname{int}(E)$ 分别表示闭包与内部。

\begin{proposition}[闭包保持凸性]\label{prop:convex-closure}
设 $X$ 为拓扑向量空间，$C\subset X$ 为凸集，则其闭包 $\overline{C}$ 仍为凸集。
\end{proposition}

\begin{proof}
任取 $x,y\in \overline{C}$ 与 $\lambda\in[0,1]$。由第 \ref{sec:1-1} 节命题~\ref{prop:closure-net}，存在取值于 $C$ 的网 $(x_\alpha)_{\alpha\in D}$ 与 $(y_\beta)_{\beta\in E}$，分别收敛到 $x$ 与 $y$。令 $D\times E$ 配上坐标逐点序构成定向集，则
\[
z_{(\alpha,\beta)}:=\lambda x_\alpha+(1-\lambda)y_\beta\in C
\]
对每个 $(\alpha,\beta)\in D\times E$ 成立。加法与数乘在拓扑向量空间中连续，因此
\[
z_{(\alpha,\beta)}\to \lambda x+(1-\lambda)y.
\]
再次应用命题~\ref{prop:closure-net}，得到 $\lambda x+(1-\lambda)y\in \overline{C}$。故 $\overline{C}$ 为凸集。
\end{proof}

\begin{proposition}[凸包与仿射包的基本性质]\label{prop:convex-hull-properties}
设 $E\subset X$，其中 $X$ 为实向量空间，则有：
\begin{enumerate}
  \item $\operatorname{conv}(E)$ 是包含 $E$ 的所有凸集的交，且
  \[
  \operatorname{conv}(E)
  =
  \left\{
    \sum_{j=1}^{m}\lambda_j x_j :
    m\in\mathbb N,\ x_j\in E,\ \lambda_j\ge 0,\ \sum_{j=1}^{m}\lambda_j=1
  \right\};
  \]
  \item $\operatorname{aff}(E)$ 是包含 $E$ 的所有仿射集的交，且
  \[
  \operatorname{aff}(E)
  =
  \left\{
    \sum_{j=1}^{m}\lambda_j x_j :
    m\in\mathbb N,\ x_j\in E,\ \lambda_j\in\mathbb R,\ \sum_{j=1}^{m}\lambda_j=1
  \right\};
  \]
  \item $\operatorname{cone}(E)$ 是包含 $E$ 的所有凸锥的交，且
  \[
  \operatorname{cone}(E)
  =
  \left\{
    \sum_{j=1}^{m}\lambda_j x_j :
    m\in\mathbb N,\ x_j\in E,\ \lambda_j\ge 0
  \right\}.
  \]
\end{enumerate}
\end{proposition}

\begin{proof}
只证明第一条，其余两条完全类似。记右端集合为 $C_0$。显然 $E\subset C_0$，且两组有限凸组合再做凸组合，仍是有限凸组合，因此 $C_0$ 是凸集。于是 $\operatorname{conv}(E)\subset C_0$。

反过来，若 $C$ 是任意包含 $E$ 的凸集，则由于 $C$ 对两点凸组合闭合，归纳可得它对任意有限凸组合都闭合，所以 $C_0\subset C$。因此 $C_0$ 包含于所有包含 $E$ 的凸集之中，故 $C_0$ 就是这些集合的交，即 $\operatorname{conv}(E)$。
\end{proof}

\paragraph{无穷维中的第一个陷阱}

\begin{remark}
在有限维空间中，很多读者会不自觉地把“闭凸”与“具有明显体积”联系起来。但在无穷维里，闭凸集完全可能没有任何拓扑内点。最简单的例子是 Hilbert 空间中的真闭子空间：它是闭、凸、仿射的，但内部为空。后面第 \ref{sec:4-2} 节引入代数内部与相对内部，正是为了修补这类现象。
\end{remark}

\begin{example}[最小抽象例子：函数空间中的正锥]\label{ex:function-space-positive-cone}
设 $X=C([0,1])$ 或 $X=L^p(0,1)$，定义
\[
X_+:=\{u\in X:u(t)\ge 0 \text{ a.e.}\}.
\]
则 $X_+$ 是一个凸锥，因为非负函数的非负线性组合仍非负。把这个例子放在这里，是为了给后文的广义不等式、正算子与锥约束提供最小抽象原型：在无穷维优化中，“$x\ge 0$”往往不再是逐坐标条件，而是属于某个函数锥的包含关系。
\end{example}

\begin{example}[有限维坍缩例子：多面体与欧氏球]\label{ex:polytope-and-ball}
在 $\mathbb R^n$ 中，半空间的有限交
\[
P=\{x\in\mathbb R^n:Ax\le b\}
\]
是凸集，非负正交第一象限是凸锥，而欧氏闭球
\[
\{x\in\mathbb R^n:\|x\|_2\le r\}
\]
则是闭凸集。把这个例子放在这里，是为了保留 Boyd 中最熟悉的几何母语：多面体、范数球与锥都是本节抽象定义的有限维投影，而不是另一套平行理论。
\end{example}

\begin{example}[算法触点例子：正则化问题的水平集]\label{ex:regularization-level-set}
设 $H$ 为 Hilbert 空间，$A\in\mathcal B(H,Y)$ 为有界线性算子，考虑泛函
\[
\Phi(u)=\frac12\|Au-b\|_Y^2+\lambda\|u\|_H^2,
\qquad \lambda>0.
\]
其下水平集
\[
L_c:=\{u\in H:\Phi(u)\le c\}
\]
是闭凸集。把这个例子放在这里，是为了说明“凸集”并不只是静态几何对象，它们正是迭代算法与存在性证明实际操作的可行域。后文讨论弱紧性、投影与邻近算子时，几乎总要先确认某个水平集是否属于本节建立的框架。
\end{example}

\subsection*{有限维对照}

Boyd 第 2 章中的凸集、锥、范数球和水平集，大多画在 $\mathbb R^2$ 或 $\mathbb R^3$ 中，因此读者很容易把“可视化”误认为“定义本身”。本节强调的恰恰相反：定义~\ref{def:convex-set} 到定义~\ref{def:convex-cone} 完全不依赖维数，真正依赖空间结构的是闭包、内部与紧性。有限维里许多几何事实因为欧氏拓扑过于良性而被自动吞掉；无穷维里则必须先把对象语法写清，后面的分离、对偶和极值存在才有可靠的起点。

\subsection*{习题（待补）}

\begin{exercise}[闭包保凸的序列版桥接题]\label{exr:convex-closure-sequence-placeholder}
补一题：在度量线性空间中把命题~\ref{prop:convex-closure} 改写成序列证明，并说明它与命题~\ref{prop:closure-net} 的网证明相比，在哪些一般拓扑环境下会失效。
\end{exercise}

\begin{exercise}[凸包有限表示占位]\label{exr:convex-hull-placeholder}
补一题：证明命题~\ref{prop:convex-hull-properties} 中的凸包公式，并把它退回到 $\mathbb R^n$ 的多面体情形，解释为什么有限维中常常只需追踪有限个顶点。
\end{exercise}


\section{代数内部与相对内部}

\label{sec:4-2}

上一节定义~\ref{def:convex-set} 已经说明，凸性本身不依赖维数，但“内部”却高度依赖所选拓扑。在有限维优化里，很多资格条件都以“存在严格可行点”为表述核心，而这里的“严格”往往就是“落在内部”。进入无穷维以后，这个说法很快会失效：一个集合即使在几何上明显具有厚度，也可能因为嵌在真闭子空间中、或者因为可行方向过于受限，而在整体拓扑下没有任何内点。因此，本节要把通常内点替换为两种更稳定的对象：代数内部描述沿所有方向的小范围可行性，相对内部则描述沿仿射包的拓扑厚度。后面的 Slater 条件、严格分离与强对偶，都将依赖这一步修正。

\subsection{两种内部概念}

\begin{definition}[代数内部]\label{def:algebraic-interior}
设 $X$ 为实向量空间，$C\subset X$。称点 $x\in C$ 属于 $C$ 的\emph{代数内部}，若对任意方向 $y\in X$，都存在 $\varepsilon>0$ 使得
\[
x+ty\in C,\qquad \forall |t|<\varepsilon.
\]
所有此类点构成的集合记为 $\operatorname{core}(C)$。
\end{definition}

\begin{definition}[相对内部]\label{def:relative-interior}
设 $X$ 为拓扑向量空间，$C\subset X$。将仿射包 $\operatorname{aff}(C)$ 赋予由 $X$ 诱导的子空间拓扑，则 $C$ 在该拓扑下的内部称为 $C$ 的\emph{相对内部}，记为
\[
\operatorname{ri}(C).
\]
\end{definition}

若 $\operatorname{aff}(C)=X$，则 $\operatorname{ri}(C)=\operatorname{int}(C)$。但在一般情形下，$\operatorname{ri}(C)$ 往往比整体内部更合适，因为优化问题的自然可行域常常落在某个仿射切片中，而不是充满整个空间。

\begin{proposition}[相对内部的基本性质]\label{prop:ri-convex-properties}
设 $X$ 为拓扑向量空间，$C\subset X$ 为非空凸集，则：
\begin{enumerate}
  \item $\operatorname{ri}(C)$ 仍为凸集；
  \item 若 $x\in \operatorname{ri}(C)$，$y\in C$，且 $\lambda\in[0,1)$，则
  \[
  (1-\lambda)x+\lambda y\in \operatorname{ri}(C).
  \]
\end{enumerate}
在有限维空间中，若 $C$ 为凸集，则 $\operatorname{core}(C)=\operatorname{ri}(C)$。
\end{proposition}

\begin{proof}
只需在仿射空间 $\operatorname{aff}(C)$ 中证明。对第一条，任取 $x_1,x_2\in \operatorname{ri}(C)$ 与 $\theta\in(0,1)$。存在 $\operatorname{aff}(C)$ 中的邻域 $U_1,U_2$，使
\[
x_1+U_1\subset C,\qquad x_2+U_2\subset C.
\]
于是
\[
\theta x_1+(1-\theta)x_2+\theta U_1+(1-\theta)U_2\subset C.
\]
而 $\theta U_1+(1-\theta)U_2$ 仍是 $\operatorname{aff}(C)$ 中 $0$ 的邻域，故 $\theta x_1+(1-\theta)x_2\in \operatorname{ri}(C)$。

对第二条，取 $x\in\operatorname{ri}(C)$，则存在 $\operatorname{aff}(C)$ 中 $0$ 的邻域 $U$ 使 $x+U\subset C$。令
\[
z=(1-\lambda)x+\lambda y.
\]
则
\[
z+(1-\lambda)U=(1-\lambda)(x+U)+\lambda y\subset C,
\]
故 $z\in\operatorname{ri}(C)$。

最后说明有限维情形下 $\operatorname{core}(C)=\operatorname{ri}(C)$。在有限维空间中，$\operatorname{aff}(C)$ 自身是有限维拓扑向量空间，因此任何代数意义上的吸收邻域都等价于某个拓扑邻域。于是点在所有方向上可作小扰动，等价于它在仿射包中拥有一个真正的开邻域。故二者相等。
\end{proof}

\paragraph{为何通常内点不够}

\begin{remark}
若 $M$ 是 Hilbert 空间 $H$ 的真闭子空间，则闭单位球
\[
B_M:=\{x\in M:\|x\|_H\le 1\}
\]
作为 $H$ 的子集是闭凸集，但其内部为空；然而作为 $M$ 内的集合，它的相对内部却是
\[
\{x\in M:\|x\|_H<1\}.
\]
这说明一旦可行域天然落在某个仿射切片内，仅用整体拓扑中的 $\operatorname{int}(C)$ 描述“严格可行”就会把所有有意义的点误删掉。
\end{remark}

\subsection{core 与分离之间的桥梁}

代数内部的真正价值不在于它本身比相对内部“更一般”，而在于它能把几何厚度转写成一个可供 Hahn--Banach 定理使用的次线性泛函。这个泛函正是 Minkowski 泛函的抽象版本。

\begin{proposition}[core 与分离的桥梁]\label{prop:core-separation-bridge}
设 $X$ 为实向量空间，$C\subset X$ 为凸集，且 $0\in \operatorname{core}(C)$。定义
\[
p_C(x):=\inf\{t>0:x\in tC\},\qquad x\in \operatorname{span}(C).
\]
则：
\begin{enumerate}
  \item $C$ 在 $\operatorname{span}(C)$ 中是吸收的，因此 $p_C(x)<\infty$ 对每个 $x\in \operatorname{span}(C)$ 成立；
  \item $p_C$ 在 $\operatorname{span}(C)$ 上是次线性泛函；
  \item 若 $p_C(x)<1$，则 $x\in C$。
\end{enumerate}
\end{proposition}

\begin{proof}
由 $0\in\operatorname{core}(C)$，对任意 $x\in\operatorname{span}(C)$，存在 $\varepsilon>0$ 使
\[
tx\in C,\qquad \forall |t|<\varepsilon.
\]
于是
\[
x=\varepsilon^{-1}(\varepsilon x)\in \varepsilon^{-1}C,
\]
故 $C$ 吸收 $\operatorname{span}(C)$ 中每个向量，从而 $p_C(x)<\infty$。

其次证明齐次性。若 $\alpha>0$，则
\[
p_C(\alpha x)
=
\inf\{t>0:\alpha x\in tC\}
=
\alpha \inf\{s>0:x\in sC\}
=
\alpha p_C(x).
\]
而 $p_C(0)=0$ 显然成立。再证明次可加性。任取 $\varepsilon>0$，取 $s,t>0$ 满足
\[
x\in sC,\quad y\in tC,\quad s<p_C(x)+\varepsilon,\quad t<p_C(y)+\varepsilon.
\]
于是存在 $u,v\in C$ 使 $x=su$、$y=tv$。由于 $C$ 凸，
\[
\frac{s}{s+t}u+\frac{t}{s+t}v\in C,
\]
从而
\[
x+y=(s+t)\left(\frac{s}{s+t}u+\frac{t}{s+t}v\right)\in (s+t)C.
\]
故
\[
p_C(x+y)\le s+t<p_C(x)+p_C(y)+2\varepsilon.
\]
令 $\varepsilon\downarrow 0$ 即得
\[
p_C(x+y)\le p_C(x)+p_C(y).
\]
因此 $p_C$ 是次线性的。

最后若 $p_C(x)<1$，取 $t<1$ 使 $x\in tC$，则存在 $u\in C$ 满足 $x=tu$。由于 $0\in C$ 且 $C$ 凸，
\[
x=tu+(1-t)0\in C.
\]
\end{proof}

\begin{remark}
命题~\ref{prop:core-separation-bridge} 解释了为什么第 \ref{sec:4-3} 节中的解析型 Hahn--Banach 定理会自然介入分离理论。只要一个凸集在某个点具有代数内部，就可以从该点出发构造 Minkowski 泛函，并把“集合分离”转写为“在线性子空间上给定一个受次线性泛函支配的线性泛函，再向外延拓”。这正是无穷维版本 Slater 条件的底层机制。
\end{remark}

\begin{example}[最小抽象例子：真子空间中的单位球]\label{ex:relative-interior-subspace}
设 $M$ 是 Hilbert 空间 $H$ 的真闭子空间，考虑
\[
C=\{x\in M:\|x\|_H\le 1\}.
\]
则 $C$ 作为 $H$ 的子集内部为空，但
\[
\operatorname{ri}(C)=\{x\in M:\|x\|_H<1\}.
\]
把这个例子放在这里，是为了说明“厚度”必须沿着自然的仿射包来测量；若仍坚持在整个 $H$ 里看内点，就会把本来良好的凸集错判成完全退化。
\end{example}

\begin{example}[有限维坍缩例子：单纯形的相对内部]\label{ex:simplex-relative-interior}
在 $\mathbb R^n$ 中考虑概率单纯形
\[
\Delta_n=\left\{x\in\mathbb R^n:x_i\ge 0,\ \sum_{i=1}^{n}x_i=1\right\}.
\]
它的仿射包是超平面 $\sum_i x_i=1$，并且
\[
\operatorname{ri}(\Delta_n)=\left\{x\in\Delta_n:x_i>0,\ 1\le i\le n\right\}.
\]
把这个例子放在这里，是为了回收 Boyd 中最常见的相对内部语言：在有限维里，$\operatorname{ri}$ 正是处理等式约束与不等式严格可行性的标准工具。
\end{example}

\begin{example}[应用触点例子：半无限约束中的严格可行性]\label{ex:semi-infinite-slater}
设 $T$ 为紧空间，考虑半无限规划可行域
\[
\mathcal F=\{x\in\mathbb R^n:a(t)^\top x\le b(t),\ \forall t\in T\},
\]
其中 $a(\cdot)$ 与 $b(\cdot)$ 连续。若存在 $\bar x$ 使
\[
a(t)^\top \bar x<b(t),\qquad \forall t\in T,
\]
且不等式余量在 $T$ 上一致正，则 $\bar x$ 落在 $\mathcal F$ 的相对内部，甚至在适当线性包内落在其 core 中。把这个例子放在这里，是为了提前说明第 7 章广义 Slater 条件的几何来源：严格可行并不是一句口号，而是某种内部结构非空。
\end{example}

\subsection*{有限维对照}

Boyd 在处理相对内部、严格可行性与约束资格时，默认工作在有限维 Euclidean 空间中，因此 $\operatorname{ri}(C)$ 已足够强大，很多时候甚至可以不再区分它与更一般的代数内部。本节的修正是：一旦进入无穷维，仅靠 $\operatorname{ri}$ 仍可能不足，因为仿射包本身未必提供合适的可行方向结构。此时 $\operatorname{core}(C)$ 与命题~\ref{prop:core-separation-bridge} 中的 Minkowski 泛函，才是把严格分离与强对偶推进下去的真正桥梁。

\subsection*{习题（待补）}

\begin{exercise}[有限维中 core 与 ri 一致占位]\label{exr:core-ri-placeholder}
补一题：在有限维空间中证明命题~\ref{prop:ri-convex-properties} 的最后一句，即对任意凸集 $C$ 有 $\operatorname{core}(C)=\operatorname{ri}(C)$。
\end{exercise}

\begin{exercise}[Minkowski 泛函桥接题占位]\label{exr:minkowski-gauge-placeholder}
补一题：在命题~\ref{prop:core-separation-bridge} 的条件下，证明若 $C$ 还是平衡集，则 $p_C$ 成为 $\operatorname{span}(C)$ 上的半范数，并说明它与第 \ref{sec:4-3} 节 Hahn--Banach 延拓的关系。
\end{exercise}


\section{Hahn-Banach 定理}

\label{sec:4-3}

从有限维凸优化的图像直觉看，超平面分离似乎是一个纯几何命题；但在无穷维中，真正起作用的并不是“画一条线把两个集合隔开”，而是在线性子空间上先构造一个受控的线性泛函，再把它延拓到整个空间。这个延拓原理就是 Hahn--Banach 定理。它一方面说明连续线性泛函足够丰富，可以探测凸集的边界；另一方面又把后文的对偶变量、支撑函数和 KKT 乘子统一到同一个对象类型中。特别地，第 \ref{sec:3-3} 节定义~\ref{def:adjoint-operator} 已经告诉我们：一旦对偶变量出现，它们就要通过伴随算子回到原问题空间。因此，本节给出的分离结论并不是孤立的几何装饰，而是第 5 章次微分和第 6 章共轭理论的基础接口。

\subsection{次线性泛函与一维延拓}

\begin{definition}[次线性泛函]\label{def:sublinear-functional}
设 $X$ 为实向量空间。映射
\[
p\colon X\to \mathbb R
\]
若满足
\begin{enumerate}
  \item $p(x+y)\le p(x)+p(y)$，对任意 $x,y\in X$；
  \item $p(\lambda x)=\lambda p(x)$，对任意 $\lambda\ge 0$ 与任意 $x\in X$，
\end{enumerate}
则称 $p$ 为 $X$ 上的\emph{次线性泛函}。
\end{definition}

\begin{lemma}[一步延拓引理]\label{lem:hb-one-step}
设 $M$ 是实向量空间 $X$ 的线性子空间，$p$ 是 $X$ 上的次线性泛函，$f_0\colon M\to\mathbb R$ 为线性泛函且满足
\[
f_0(x)\le p(x),\qquad \forall x\in M.
\]
若 $x_0\in X\setminus M$，则存在定义在 $M+\mathbb R x_0$ 上的线性泛函 $\widetilde f$，使
\[
\widetilde f|_M=f_0,\qquad \widetilde f(x)\le p(x),\ \forall x\in M+\mathbb R x_0.
\]
\end{lemma}

\begin{proof}
我们寻找常数 $c\in\mathbb R$，并定义
\[
\widetilde f(m+t x_0):=f_0(m)+tc,\qquad m\in M,\ t\in\mathbb R.
\]
要使 $\widetilde f\le p$，等价于要求
\[
f_0(m)+tc\le p(m+t x_0),\qquad \forall m\in M,\ \forall t\in\mathbb R.
\]
只需考察 $t=1$ 与 $t=-1$，即可得到约束
\[
f_0(m)+c\le p(m+x_0),\qquad
f_0(m)-c\le p(m-x_0),\qquad \forall m\in M.
\]
亦即
\[
f_0(m)-p(m-x_0)\le c\le p(m+x_0)-f_0(m),\qquad \forall m\in M.
\]
因此只要证明区间
\[
\left[
\sup_{m\in M}\bigl(f_0(m)-p(m-x_0)\bigr),\
\inf_{m\in M}\bigl(p(m+x_0)-f_0(m)\bigr)
\right]
\]
非空即可。

任取 $m_1,m_2\in M$。由 $f_0\le p$ 及次线性，
\[
f_0(m_1)+f_0(m_2)
=f_0(m_1+m_2)
\le p(m_1+m_2)
=p((m_1-x_0)+(m_2+x_0))
\le p(m_1-x_0)+p(m_2+x_0).
\]
移项得
\[
f_0(m_1)-p(m_1-x_0)\le p(m_2+x_0)-f_0(m_2).
\]
对左边取上确界、右边取下确界，便知该区间非空。任选其中一点 $c$，即可得到所需的 $\widetilde f$。
\end{proof}

\begin{theorem}[Hahn--Banach 的解析形式]\label{thm:hahn-banach-analytic}
设 $X$ 为实向量空间，$M\subset X$ 为线性子空间，$p$ 为 $X$ 上的次线性泛函，$f_0\colon M\to\mathbb R$ 为线性泛函且满足
\[
f_0(x)\le p(x),\qquad \forall x\in M.
\]
则存在定义在整个 $X$ 上的线性泛函 $f\colon X\to\mathbb R$，使
\[
f|_M=f_0,\qquad f(x)\le p(x),\ \forall x\in X.
\]
\end{theorem}

\begin{proof}
考虑所有满足下列条件的对 $(N,g)$ 所成集合：
\begin{enumerate}
  \item $M\subset N\subset X$，且 $N$ 是线性子空间；
  \item $g\colon N\to\mathbb R$ 线性；
  \item $g|_M=f_0$；
  \item $g(x)\le p(x)$ 对任意 $x\in N$ 成立。
\end{enumerate}
以“延拓”定义偏序：若 $(N_1,g_1)\preceq (N_2,g_2)$，则指 $N_1\subset N_2$ 且 $g_2|_{N_1}=g_1$。任意全序链的并仍给出一个满足上述条件的上界，因此由 Zorn 引理，存在极大元 $(N^\ast,g^\ast)$。

若 $N^\ast\neq X$，取 $x_0\in X\setminus N^\ast$。由引理~\ref{lem:hb-one-step}，$g^\ast$ 可以延拓到 $N^\ast+\mathbb R x_0$ 上且仍被 $p$ 支配，这与极大性矛盾。故 $N^\ast=X$。令 $f=g^\ast$，即得所求。
\end{proof}

\begin{remark}
若 $X$ 还是赋范空间，而次线性泛函 $p$ 满足
\[
p(x)\le C\|x\|,\qquad \forall x\in X,
\]
则定理~\ref{thm:hahn-banach-analytic} 给出的延拓 $f$ 自动连续，并满足 $\|f\|\le C$。因此，解析形式的 Hahn--Banach 定理不仅保证“存在足够多的线性泛函”，还保证在 Banach 空间里它们属于连续对偶空间 $X^\ast$，这正是后文分离与对偶所需的对象层级。
\end{remark}

\subsection{几何分离}

\begin{theorem}[Hahn--Banach 的几何形式]\label{thm:hahn-banach-geometric}
设 $X$ 为实赋范空间，$C\subset X$ 为非空开凸集，$x_0\notin C$。则存在非零连续线性泛函 $f\in X^\ast$ 与常数 $\alpha\in\mathbb R$，使得
\[
f(x)<\alpha\le f(x_0),\qquad \forall x\in C.
\]
\end{theorem}

\begin{proof}
任取 $a\in C$，令
\[
U:=C-a,\qquad z_0:=x_0-a.
\]
则 $U$ 是包含原点的开凸集，而 $z_0\notin U$。定义 $U$ 的 Minkowski 泛函
\[
p_U(x):=\inf\{t>0:x\in tU\},\qquad x\in X.
\]
由于 $U$ 开、凸且 $0\in U$，由命题~\ref{prop:core-separation-bridge} 的思想可知 $p_U$ 是次线性的。又因为存在 $r>0$ 使
\[
B(0,r)\subset U,
\]
故对任意 $x\in X$ 都有
\[
p_U(x)\le \frac{\|x\|}{r},
\]
于是 $p_U$ 还被一个连续半范数控制。

由于 $z_0\notin U$，故
\[
p_U(z_0)\ge 1.
\]
在一维子空间 $M:=\mathbb R z_0$ 上定义线性泛函
\[
f_0(tz_0):=t.
\]
我们验证 $f_0\le p_U$。若 $t\ge 0$，则
\[
f_0(tz_0)=t\le t\,p_U(z_0)=p_U(tz_0).
\]
若 $t<0$，则 $f_0(tz_0)=t\le 0\le p_U(tz_0)$。于是由定理~\ref{thm:hahn-banach-analytic}，存在线性泛函 $f$ 延拓 $f_0$ 且满足
\[
f(x)\le p_U(x),\qquad \forall x\in X.
\]
再由 $p_U(x)\le \|x\|/r$ 知 $f\in X^\ast$。并且
\[
f(z_0)=f_0(z_0)=1,
\]
故 $f\neq 0$。

任取 $x\in C$，则 $x-a\in U$。由于 $U$ 是开集且包含原点，存在 $\eta>0$ 使
\[
(1+\eta)(x-a)\in U.
\]
于是
\[
x-a\in \frac{1}{1+\eta}U,
\]
从而
\[
p_U(x-a)\le \frac{1}{1+\eta}<1.
\]
故
\[
f(x)-f(a)=f(x-a)\le p_U(x-a)<1=f(z_0)=f(x_0)-f(a),
\]
即
\[
f(x)<f(x_0),\qquad \forall x\in C.
\]
取 $\alpha:=f(x_0)$ 即得结论。
\end{proof}

\begin{theorem}[严格分离定理]\label{thm:strict-separation}
设 $X$ 为实赋范空间，$C\subset X$ 为非空闭凸集，$x_0\notin C$。则存在非零连续线性泛函 $f\in X^\ast$ 与实数 $\alpha<\beta$，使
\[
f(x)\le \alpha<\beta\le f(x_0),\qquad \forall x\in C.
\]
\end{theorem}

\begin{proof}
记
\[
d:=\operatorname{dist}(x_0,C)>0.
\]
令
\[
V:=C+B\left(0,\frac d2\right).
\]
则 $V$ 是非空开凸集，且 $x_0\notin V$。由定理~\ref{thm:hahn-banach-geometric}，存在非零 $f\in X^\ast$ 与 $\gamma\in\mathbb R$，使
\[
f(v)<\gamma\le f(x_0),\qquad \forall v\in V.
\]
由于 $f\neq 0$，可取某个 $z\in X$ 使 $f(z)>0$，再缩放得到 $h\in B(0,d/2)$ 且 $f(h)>0$。于是对任意 $x\in C$，都有
\[
x+h\in V,
\]
从而
\[
f(x)+f(h)=f(x+h)<\gamma.
\]
因此
\[
f(x)<\gamma-f(h),\qquad \forall x\in C.
\]
取
\[
\alpha:=\gamma-f(h),\qquad \beta:=\gamma,
\]
便得到
\[
f(x)\le \alpha<\beta\le f(x_0),\qquad \forall x\in C.
\]
\end{proof}

\begin{proposition}[支撑超平面]\label{prop:supporting-hyperplane}
设 $X$ 为实赋范空间，$C\subset X$ 为非空闭凸集且 $\operatorname{int}(C)\neq\varnothing$。若 $x_0\in \partial C$，则存在非零连续线性泛函 $f\in X^\ast$ 使
\[
f(x)\le f(x_0),\qquad \forall x\in C.
\]
因此超平面
\[
H:=\{x\in X:f(x)=f(x_0)\}
\]
称为 $C$ 在 $x_0$ 处的一个\emph{支撑超平面}。
\end{proposition}

\begin{proof}
对开凸集 $\operatorname{int}(C)$ 与外点 $x_0\notin \operatorname{int}(C)$ 应用定理~\ref{thm:hahn-banach-geometric}，可得非零 $f\in X^\ast$ 与常数 $\alpha$，使
\[
f(y)<\alpha\le f(x_0),\qquad \forall y\in \operatorname{int}(C).
\]
任取 $u\in \operatorname{int}(C)$。对每个 $x\in C$ 及每个 $t\in(0,1]$，由凸性知
\[
(1-t)x+tu\in \operatorname{int}(C).
\]
于是
\[
f((1-t)x+tu)<\alpha.
\]
令 $t\downarrow 0$，由连续性得到
\[
f(x)\le \alpha,\qquad \forall x\in C.
\]
另一方面，对边界点 $x_0$ 仍有
\[
(1-t)x_0+tu\in \operatorname{int}(C),\qquad t\in(0,1],
\]
故
\[
f((1-t)x_0+tu)<\alpha.
\]
令 $t\downarrow 0$ 得
\[
f(x_0)\le \alpha.
\]
与 $\alpha\le f(x_0)$ 合并，即得
\[
f(x_0)=\alpha,\qquad f(x)\le f(x_0)\ \forall x\in C.
\]
\end{proof}

\begin{remark}
命题~\ref{prop:supporting-hyperplane} 中出现的 $f\in X^\ast$，就是后文次梯度与对偶变量的原型。对约束问题
\[
Tu=b
\]
而言，乘子首先属于 $Y^\ast$，随后通过第 \ref{sec:3-3} 节定义~\ref{def:adjoint-operator} 里的伴随算子 $T^\ast$ 拉回到 $X^\ast$。因此，几何分离定理与伴随算子并不是两块独立内容，而是同一套对偶语言的两端。
\end{remark}

\begin{example}[最小抽象例子：从常数函数子空间的延拓]\label{ex:hahn-banach-extension}
设 $X=C([0,1])$，$M=\operatorname{span}\{\mathbf 1\}$，定义
\[
f_0(c\mathbf 1)=c,\qquad p(u)=\|u\|_\infty.
\]
显然 $f_0\le p$。由定理~\ref{thm:hahn-banach-analytic}，$f_0$ 可以延拓为某个 $F\in X^\ast$，并满足 $\|F\|\le 1$。再由第 \ref{sec:2-3} 节定理~\ref{thm:riesz-representation}，$F$ 对应于一个有限符号 Radon 测度。把这个例子放在这里，是为了说明无穷维中的“法向量”常常并不是一个向量，而是一类测度型泛函。
\end{example}

\begin{example}[有限维坍缩例子：欧氏球的支撑超平面]\label{ex:euclidean-supporting-hyperplane}
在 $\mathbb R^n$ 中，单位闭球
\[
B=\{x\in\mathbb R^n:\|x\|_2\le 1\}
\]
在边界点 $x_0$ 处的支撑超平面由
\[
\langle x_0,x\rangle=1
\]
给出。把这个例子放在这里，是为了把 Boyd 中最熟悉的“法向量接触球面”图像精确识别为命题~\ref{prop:supporting-hyperplane} 在 Hilbert 空间里的有限维坍缩。
\end{example}

\begin{example}[应用触点例子：对偶乘子的几何来源]\label{ex:dual-multiplier-separation}
考虑凸约束问题
\[
\min_{u\in X}\ \phi(u)\qquad \text{s.t.}\qquad Tu=b,
\]
其中 $T\in\mathcal B(X,Y)$。在最优点附近，把目标的下降方向与线性约束切空间分离后，得到的正是某个 $\lambda\in Y^\ast$；再经由 $T^\ast\lambda\in X^\ast$ 把它带回原变量空间。把这个例子放在这里，是为了强调：Boyd 书里写作 $A^\top\lambda$ 的对象，在无穷维中真正依赖的是 Hahn--Banach 分离加上伴随算子，而不是矩阵技巧本身。
\end{example}

\subsection*{有限维对照}

Boyd 在第 2 章和第 5 章中大量使用超平面分离、支撑函数和对偶变量，但这些对象在有限维里往往直接以向量和矩阵写出，因此它们的泛函本性被隐藏了。本节的结论说明：有限维图像并不是另一套理论，而是定理~\ref{thm:hahn-banach-geometric}、定理~\ref{thm:strict-separation} 与命题~\ref{prop:supporting-hyperplane} 在 Euclidean 空间中的压缩版本。真正变化的只是表面符号；底层机制始终是“线性泛函足够丰富，可以把凸集与点、集合与集合、原问题与对偶问题分开”。

\subsection*{习题（待补）}

\begin{exercise}[一步延拓公式占位]\label{exr:hahn-banach-one-step-placeholder}
补一题：完整推导引理~\ref{lem:hb-one-step} 中常数 $c$ 所落入的区间，并说明为何该区间非空正是次线性条件的体现。
\end{exercise}

\begin{exercise}[支撑超平面桥接题占位]\label{exr:supporting-hyperplane-placeholder}
补一题：在 Hilbert 空间中把命题~\ref{prop:supporting-hyperplane} 退化到闭球情形，直接计算支撑超平面的法向量，并与例~\ref{ex:euclidean-supporting-hyperplane} 对照。
\end{exercise}


\section{极点与 Krein-Milman 定理}

\label{sec:4-4}

前面三节说明了如何用线性泛函描述凸集的边界，但“边界”本身仍然太大。对优化而言，真正决定结构的往往不是整个边界，而是其中不能再被非平凡凸组合分解的那一部分。这便引出极点。在线性规划里，极点退化为顶点；在稀疏优化里，极点对应原子；在概率测度空间里，极点对应 Dirac 质量。Krein-Milman 定理的意义在于：它把这些“极端元”提升为一般紧凸集的生成机制，说明一个紧凸集可以由其极点通过闭凸包完全恢复。这样一来，第 \ref{sec:4-3} 节的分离理论不再只告诉我们“边界可被支撑”，还告诉我们“结构最终集中在边界最尖锐的部分”。

\subsection{极点与面}

\begin{definition}[极点]\label{def:extreme-point}
设 $C$ 为实向量空间中的凸集。若 $x\in C$ 满足：只要
\[
x=\lambda y+(1-\lambda)z,\qquad y,z\in C,\ \lambda\in(0,1),
\]
就必有 $y=z=x$，则称 $x$ 为 $C$ 的一个\emph{极点}。所有极点组成的集合记为
\[
\operatorname{ext}(C).
\]
\end{definition}

为了组织极点，常引入\emph{面}的概念：若 $F\subset C$ 为非空凸子集，并且只要
\[
\lambda y+(1-\lambda)z\in F,\qquad y,z\in C,\ \lambda\in(0,1),
\]
就推出 $y,z\in F$，则称 $F$ 为 $C$ 的一个面。单点面恰好对应极点；换言之，$x$ 是极点，当且仅当 $\{x\}$ 是一个面。

\begin{lemma}[暴露面]\label{lem:exposed-face}
设 $X$ 为 Hausdorff 局部凸空间，$K\subset X$ 为非空紧凸集，$f\in X^\ast$ 为连续线性泛函。令
\[
M:=\max_{x\in K} f(x),\qquad F:=\{x\in K:f(x)=M\}.
\]
则 $F$ 是 $K$ 的一个非空紧面。
\end{lemma}

\begin{proof}
由于 $K$ 紧且 $f$ 连续，$f$ 在 $K$ 上取到最大值，所以 $F$ 非空。又因为 $F$ 是闭集与紧集 $K$ 的交，所以 $F$ 紧。对任意 $x_1,x_2\in F$ 与 $\lambda\in[0,1]$，
\[
f(\lambda x_1+(1-\lambda)x_2)=\lambda M+(1-\lambda)M=M,
\]
故 $F$ 凸。

若 $\lambda y+(1-\lambda)z\in F$，其中 $y,z\in K$ 且 $\lambda\in(0,1)$，则
\[
M=f(\lambda y+(1-\lambda)z)=\lambda f(y)+(1-\lambda)f(z)\le \lambda M+(1-\lambda)M=M.
\]
因此必须有 $f(y)=f(z)=M$，即 $y,z\in F$。故 $F$ 为面。
\end{proof}

\begin{theorem}[Krein--Milman 定理]\label{thm:krein-milman}
设 $X$ 为 Hausdorff 局部凸空间，$K\subset X$ 为非空紧凸集。则
\[
K=\overline{\operatorname{conv}(\operatorname{ext}(K))}.
\]
\end{theorem}

\begin{proof}
证明分两步进行。

\paragraph{第一步：$K$ 至少有一个极点}

考虑 $K$ 的所有非空紧面所成集合，并以反向包含关系排序。任一全序链的交仍是非空紧面：非空性来自紧集族有限交性质，面性质则由定义直接继承。由 Zorn 引理，存在一个极小非空紧面 $F$。

若 $F$ 不是单点集，则取互异点 $x_1,x_2\in F$。由于 $X$ 局部凸且 Hausdorff，连续线性泛函足够分离点，因此存在 $f\in X^\ast$ 使
\[
f(x_1)\neq f(x_2).
\]
由引理~\ref{lem:exposed-face}，集合
\[
F_1:=\left\{x\in F:f(x)=\max_{y\in F} f(y)\right\}
\]
是 $F$ 的非空紧面。由于 $f$ 在 $F$ 上非常数，$F_1$ 是 $F$ 的真子集，这与 $F$ 的极小性矛盾。故 $F$ 必为单点集 $\{x\}$。此时 $\{x\}$ 为 $K$ 的面，所以 $x$ 是 $K$ 的极点。

\paragraph{第二步：极点的闭凸包已经生成整个 $K$}

记
\[
C:=\overline{\operatorname{conv}(\operatorname{ext}(K))}.
\]
显然 $C\subset K$，且 $C$ 为闭凸集。若 $C\neq K$，取 $x_0\in K\setminus C$。由于 $K$ 所在空间局部凸，可对点与闭凸集应用 Hahn--Banach 型分离；这正是定理~\ref{thm:hahn-banach-geometric} 的局部凸推广。于是存在连续线性泛函 $f\in X^\ast$ 与常数 $\alpha<\beta$，使
\[
f(x)\le \alpha<\beta\le f(x_0),\qquad \forall x\in C.
\]

令
\[
F:=\left\{x\in K:f(x)=\max_{y\in K} f(y)\right\}.
\]
由引理~\ref{lem:exposed-face}，$F$ 是 $K$ 的非空紧面。由第一步，$F$ 至少包含一个极点 $u$。而面上的极点仍是整个集合的极点：若
\[
u=\lambda y+(1-\lambda)z,\qquad y,z\in K,\ \lambda\in(0,1),
\]
则因 $u\in F$ 且 $F$ 为面，必有 $y,z\in F$；再由 $u$ 在 $F$ 中是极点，得 $y=z=u$。于是
\[
u\in \operatorname{ext}(K)\subset C.
\]
但另一方面，$u\in F$ 且 $x_0\in K$，所以
\[
f(u)=\max_{y\in K}f(y)\ge f(x_0)\ge \beta>\alpha.
\]
这与 $u\in C$ 时的 $f(u)\le \alpha$ 矛盾。故 $C=K$，定理得证。
\end{proof}

\begin{corollary}[有限维多面体的极点描述]\label{cor:finite-dimensional-polytope-extreme}
设 $P\subset\mathbb R^n$ 为紧多面体，则
\[
P=\operatorname{conv}(\operatorname{ext}(P)).
\]
特别地，在有限维情形中，多面体的顶点恰好就是它的极点。
\end{corollary}

\begin{proof}
紧多面体是有限维空间中的紧凸集，因此可直接应用定理~\ref{thm:krein-milman}。由于有限维中闭包运算对有限凸包不再引入额外复杂性，得到
\[
P=\operatorname{conv}(\operatorname{ext}(P)).
\]
而顶点正是不能写成其它可行点非平凡凸组合的点，所以与定义~\ref{def:extreme-point} 完全一致。
\end{proof}

\paragraph{结构意义}

\begin{remark}
定理~\ref{thm:krein-milman} 把第 \ref{sec:4-3} 节的分离原理推进了一步。分离只告诉我们：若一点不在某个闭凸集中，则存在连续线性泛函把它隔开；而 Krein-Milman 进一步指出：只要集合还具有紧性，这些支撑泛函最终会把结构压缩到极点上。后文无论是讨论原子范数、测度型最优策略还是线性规划顶点，都是这一思想在不同空间中的再现。
\end{remark}

\begin{example}[最小抽象例子：概率测度的 Dirac 极点]\label{ex:dirac-extreme-point}
设 $K$ 为紧 Hausdorff 空间，$\mathcal P(K)$ 表示其上所有 Borel 概率测度的集合。借助第 \ref{sec:2-3} 节定理~\ref{thm:riesz-representation}，可把 $\mathcal P(K)$ 看成 $C(K)^\ast$ 中的一个弱$^\ast$紧凸集。对任意 $x\in K$，Dirac 测度 $\delta_x$ 都是 $\mathcal P(K)$ 的极点：若
\[
\delta_x=\lambda\mu_1+(1-\lambda)\mu_2,\qquad \lambda\in(0,1),
\]
则对任意不含 $x$ 的闭集 $F$ 有 $\delta_x(F)=0$，从而 $\mu_1(F)=\mu_2(F)=0$；再利用总质量为 $1$ 可知 $\mu_1=\mu_2=\delta_x$。把这个例子放在这里，是为了说明无穷维随机优化中“极端策略”往往对应于原子测度。
\end{example}

\begin{example}[有限维坍缩例子：线性规划顶点]\label{ex:lp-vertex-extreme-point}
在线性规划可行域
\[
P=\{x\in\mathbb R^n:Ax\le b\}
\]
中，顶点恰好就是极点。把这个例子放在这里，是为了与 Boyd 的 LP 几何直接对齐：有限维教材里“最优解常在顶点处取得”的说法，正是推论~\ref{cor:finite-dimensional-polytope-extreme} 的线性规划版本。
\end{example}

\begin{example}[应用触点例子：$\ell^1$ 球的原子结构]\label{ex:l1-ball-extreme-points}
在 $\mathbb R^n$ 中，
\[
B_1=\{x\in\mathbb R^n:\|x\|_1\le 1\}
\]
的极点正是 $\{\pm e_1,\dots,\pm e_n\}$。把这个例子放在这里，是为了连接稀疏优化中的原子范数直觉：一个点若位于 $\ell^1$ 球的极点上，就只激活单个坐标；而一般点则是这些原子的凸组合。无穷维原子表示的语言，本质上就是把这种有限维现象推广到更大的函数或测度空间。
\end{example}

\subsection*{有限维对照}

Boyd 书中关于顶点、极点、稀疏结构和 LP 几何的大部分直觉，都可以看作定理~\ref{thm:krein-milman} 在有限维空间中的压缩版。有限维里，多面体顶点可直接画出来，$\ell^1$ 球的尖点也能肉眼识别；无穷维里则必须借助极点与闭凸包的语言来组织这些结构。换句话说，Boyd 所讨论的顶点现象并不是偶然的多面体事实，而是一般紧凸集边界理论在有限维中的最易见表现。

\subsection*{习题（待补）}

\begin{exercise}[极点与面占位]\label{exr:extreme-point-face-placeholder}
补一题：证明“$\{x\}$ 是凸集 $C$ 的面”与“$x$ 是 $C$ 的极点”等价，并说明这一命题在定理~\ref{thm:krein-milman} 的证明中起到什么作用。
\end{exercise}

\begin{exercise}[Dirac 测度桥接题占位]\label{exr:dirac-measure-placeholder}
补一题：在紧空间 $K$ 上验证例~\ref{ex:dirac-extreme-point} 中 Dirac 测度的极点性质，并进一步讨论哪些条件下 $\mathcal P(K)$ 的全部极点都由 Dirac 测度给出。
\end{exercise}

